\section{Formal potentials for a single edge}
\label{sec:edge-germs}

Let $Y$ be a compact connected oriented smooth three-manifold with
connected boundary $T\cong S^1\times S^1$.  Fix based coordinate circles
$\alpha,\beta$ on $T$.  We assume that the image of
$H^1(Y;\mathbb C)\to H^1(T;\mathbb C)$ is the line evaluating on
$\alpha$ and vanishing on $\beta$.  These are the boundary data furnished
by the edge pieces.  All cohomology in this section has complex
coefficients.  Put $\g=\mathfrak{sl}_2(\mathbb C)$, with its invariant
nondegenerate pairing $(X,Y)\mapsto\tr(XY)$.
Our goal is a formal potential whose actual gradient ideal describes
based flat systems with one peripheral holonomy fixed. We also need
the first-order change in this potential when that holonomy varies.

Choose closed one-forms $\zeta,\xi$ on $T$ with periods
\[
 \int_\alpha\zeta=\int_\beta\xi=1,
 \qquad
 \int_\beta\zeta=\int_\alpha\xi=0.
\]
They may be taken to be bump densities in the two circle coordinates.
Their supports are thin strips; choose their common basepoint outside
both strips, and choose the coordinate circles through that basepoint.
Thus $\xi$ vanishes along $\alpha$, and $\zeta$ vanishes along $\beta$.
Write
\[
 \epsilon=\int_T\zeta\wedge\xi\in\{1,-1\},
\]
where $T$ has its boundary orientation.  Choose a two-form $\eta$ of
integral one supported in a small disk disjoint from the strips.  The
restriction hypothesis gives a closed extension
$\widehat\zeta\in\Omega^1(Y)$ of $\zeta$: first choose a closed
representative with the required boundary cohomology class, and then
correct its boundary trace by the differential of an extended function.

\subsection{The boundary complex and its pairing}

Define a nonunital commutative differential graded algebra $C$ by
\begin{equation}
\label{eq:edge-boundary-complex}
 C^j=\left\{a\in\Omega^j(Y):
 a|_T\in
 \begin{cases}
  0,&j=0,\\
  \mathbb C\xi,&j=1,\\
  \mathbb C\eta,&j=2,\\
  0,&j\geq3.
 \end{cases}\right\}.
\end{equation}
In degrees at least three the boundary condition is automatic.
Multiplication is closed because degree-zero boundary values vanish,
$\xi\wedge\xi=0$, and all other positive-degree products in question
vanish on the two-dimensional boundary.

\begin{lemma}
\label{lem:edge-cohomology}
There are natural isomorphisms
\[
 H^1(Y,T)\cong H^1(C),\qquad H^2(C)\cong H^2(Y),
\]
and $H^0(C)=H^3(C)=0$.  Both middle groups have dimension
$n=b_1(Y)-1=b_2(Y)$.  Every closed one-form in $C$ has zero actual
boundary pullback.  Integration gives a perfect pairing
\[
 H^1(C)\otimes H^2(C)\longrightarrow\mathbb C.
\]
\end{lemma}

\begin{proof}
The quotient of $C$ by the relative de Rham complex has just the two
terms $\mathbb C\xi$ in degree one and $\mathbb C\eta$ in degree two,
with zero differential.  The connecting map
\[
 \mathbb C\xi\longrightarrow H^2(Y,T)
\]
is injective, since $[\xi]$ does not lie in the image of $H^1(Y)$.
The connecting map
\[
 \mathbb C\eta\longrightarrow H^3(Y,T)
\]
is an isomorphism: under the relative orientation it is boundary
integration.  The long exact sequence consequently identifies $H^1(C)$
with $H^1(Y,T)$ and $H^2(C)$ with the quotient of $H^2(Y,T)$ by the
connecting $\xi$-line.  The ordinary relative sequence identifies this
last quotient with $H^2(Y)$.  It also gives the asserted vanishing in
degrees zero and three.  Poincar\'e--Lefschetz duality gives
$\dim H^1(Y,T)=b_2(Y)$, and
$\chi(Y)=\tfrac12\chi(T)=0$ gives $b_2(Y)=b_1(Y)-1$.

If a closed one-form in $C$ has trace $v\xi$, its boundary cohomology
class is both a multiple of $[\xi]$ and an element of the $[\zeta]$-line.
Hence $v=0$.  Under the displayed isomorphisms, the integration pairing
is the relative/absolute Poincar\'e--Lefschetz pairing in degrees one
and two, so it is perfect.
\end{proof}

We keep $C$ for the scalar complex.  The tensor product $C\otimes\g$
is a dg Lie algebra, with bracket given by the wedge product combined
with the Lie bracket.  For homogeneous $a,b\in C\otimes\g$ whose
degrees sum to three, put
\[
 \langle a,b\rangle=\int_Y\tr(a\wedge b).
\]
This pairing satisfies cyclic Stokes.  Indeed, when the degrees of
$a,b$ sum to two, their boundary product is zero: either one has degree
zero or both boundary traces are multiples of $\xi$.  Thus
\begin{equation}
\label{eq:edge-cyclic-stokes}
 \langle da,b\rangle+(-1)^{|a|}\langle a,db\rangle=0.
\end{equation}
Invariance of the matrix trace gives the corresponding cyclic identity
for the bracket.  No assertion that integration of a single
degree-two element commutes with the differential is needed; its
boundary trace may be a multiple of $\eta$.

Choose representatives $i:H^*(C)\to C$, a projection
$\operatorname{pr}:C\to H^*(C)$, and a degree-minus-one map $K$ such that
\begin{equation}
\label{eq:edge-contraction}
 dK+Kd=\id-i\operatorname{pr},\qquad
 K^2=Ki=\operatorname{pr}K=0,\qquad
 \operatorname{pr}i=\id.
\end{equation}
These maps exist by splitting the vector spaces into boundaries,
cohomology representatives, and complements on which the differential
is an isomorphism onto boundaries.  Extend them by the identity on
$\g$.  In particular, the degree-one representatives have zero
boundary pullback.  Since $H^0(C)=0$, we have $Kd=\id$ on $C^0$.

All series below are formal.  Applying one of these linear maps means
applying it separately to each coefficient, which is an actual smooth
form.  No boundedness or analytic convergence of $K$ is required.

\subsection{The central slice and its potential}

Set
\[
 H=H^1(C)\otimes\g,\qquad N=\dim H=3n,
 \qquad S=\mathbb C[[h_1,\ldots,h_N]],
\]
and let $h$ denote the universal element of $H\otimes S$.  Its
coordinates generate the maximal ideal $\m$ of $S$.  The recursion
\begin{equation}
\label{eq:edge-central-recursion}
 s=i(h)-\tfrac12K[s,s]
\end{equation}
has a unique solution in $C^1\otimes\g\widehat\otimes\m$: its linear
term is $i(h)$, and each higher homogeneous term is determined by
previous terms.  It satisfies
\[
 \operatorname{pr}s=h,\qquad Ks=0,
 \qquad K P(s)=0,
 \qquad P(a)=da+\tfrac12[a,a].
\]
Conversely these three slice conditions imply
\eqref{eq:edge-central-recursion}, by
\eqref{eq:edge-contraction}.  Write $s|_T=v(h)\xi$.

We use the Chern--Simons transgression functional
\cite[Section~3]{ChernSimons1974}. For a Lie-algebra-valued one-form $a$,
define
\begin{equation}
\label{eq:edge-cs}
 \mathcal C(a)=\int_Y\tr\left(
        \tfrac12a\wedge da+\tfrac16a\wedge[a,a]\right),
 \qquad f(h)=\mathcal C(s(h)).
\end{equation}
For any coefficient variation $\dot a$, integration by parts gives
\begin{equation}
\label{eq:edge-variation}
 \dot{\mathcal C}(a)
   =\int_Y\tr(\dot a\wedge P(a))
       -\tfrac12\int_T\tr(a\wedge\dot a).
\end{equation}
In the central slice the boundary term is zero.

For a formal form $P$, its \emph{curvature ideal} means the ideal
generated by all its scalar coefficient functions, equivalently by
applying arbitrary linear functionals on the vector space of
Lie-algebra-valued forms coefficientwise.  The next calculation in
particular shows that the ideals occurring here are finitely generated.

\begin{lemma}
\label{lem:edge-central-ideal}
The curvature ideal of $s(h)$ is the Jacobian ideal
\[
 J=\Jac(f)=\left(\frac{\partial f}{\partial h_1},\ldots,
                 \frac{\partial f}{\partial h_N}\right).
\]
Moreover $f\in\m^3$.
\end{lemma}

\begin{proof}
Put $P=P(s)$ and $r=\operatorname{pr}P\in H^2(C)\otimes\g
\widehat\otimes S$.  The boundary trace of $P$ is zero, so $P$ lies coefficientwise in $C^2\otimes\g$.
Bianchi and the projected slice equation give
\begin{equation}
\label{eq:edge-central-curvature}
 dP=-[s,P],\qquad KP=0,
 \qquad P=i(r)-K[s,P].
\end{equation}
The operator $K\ad_s$ raises maximal-ideal order.  Therefore
\[
 P=(1+K\ad_s)^{-1}i(r).
\]
Every curvature coefficient lies in the ideal of the finitely many
components of $r$, and the converse holds by applying
$\operatorname{pr}$.  This is an equality of ideals, without passage to
their radicals.

By \eqref{eq:edge-variation}, the gradient entries are
\[
 \frac{\partial f}{\partial h_j}
    =\left\langle\frac{\partial s}{\partial h_j},P\right\rangle.
\]
Consequently the gradient is a square matrix of formal functions times
the vector $r$.  Its constant matrix is the perfect pairing between
$H^1(C)\otimes\g$ and $H^2(C)\otimes\g$ from
Lemma~\ref{lem:edge-cohomology}; it is invertible.  This proves the ideal
equality.  Finally $s=i(h)+O(\m^2)$ and $di(h)=0$.  Both terms in
\eqref{eq:edge-cs} therefore have order at least three.
\end{proof}

\subsection{A prescribed first-order boundary holonomy}

The central slice fixes the surviving peripheral holonomy to the
identity. The vertex construction will vary that holonomy by a
square-zero parameter. Its curvature may then have a boundary
component, so the next calculation includes that component in the
gradient ideal.

Fix $T_0\in\g$ and a scalar parameter $\tau$ with $\tau^2=0$.  Put
\[
 S_\tau=S[\tau]/(\tau^2),\qquad u=\tau T_0.
\]
The formal topology is the one at $(h,\tau)=0$.  Choose any linear
extension
\[
 K_{\mathrm{ext}}:\Omega^2(Y)\longrightarrow
       \im\bigl(K:C^2\to C^1\bigr)
\]
of the degree-two part of $K$, and tensor it with $\g$.  Define
$s=s(h,\tau)$ by
\begin{equation}
\label{eq:edge-affine-recursion}
 s=i(h)+u\widehat\zeta-\tfrac12K_{\mathrm{ext}}[s,s].
\end{equation}
The recursion again determines a unique formal series.  It reduces to
the central slice when $\tau=0$, and
\[
 s-u\widehat\zeta\in C^1\otimes\g\widehat\otimes S_\tau,
 \quad \operatorname{pr}(s-u\widehat\zeta)=h,
 \quad K(s-u\widehat\zeta)=0.
\]
Its boundary value has the form
\[
 s|_T=u\zeta+v(h,\tau)\xi.
\]
The recursion implies $K_{\mathrm{ext}}P(s)=0$.  To see this directly,
write its correction as an element $\ell$ of the original one-form
image of $K$.  The contraction identities give $Kd\ell=\ell$;
also $di(h)=d\widehat\zeta=0$.  Applying $K_{\mathrm{ext}}$ to the
curvature then gives exactly the recursion.

The boundary curvature is generally nonzero off the critical scheme:
\begin{equation}
\label{eq:edge-boundary-curvature}
 P(s)|_T=c\,\zeta\wedge\xi,\qquad c=[u,v].
\end{equation}
In particular it need not satisfy the degree-two boundary condition
defining $C$.  The following argument keeps this additional component.

\begin{lemma}
\label{lem:edge-affine-ideal}
The function
\begin{equation}
\label{eq:edge-adjusted-potential}
 f_\tau(h)=\mathcal C(s(h,\tau))
                   +\tfrac{\epsilon}{2}\tr(u v(h,\tau))
\end{equation}
has internal Jacobian ideal equal to the full curvature ideal of the
affine slice.  Here the internal derivatives are only the derivatives
with respect to the $h_j$.
\end{lemma}

\begin{proof}
Choose a scalar two-form $\rho$ on $Y$ with boundary trace
$\zeta\wedge\xi$, and put
\[
 P=P(s),\qquad P_0=P-c\rho,\qquad
 r=\operatorname{pr}P_0.
\]
Now $P_0$ has zero boundary trace and lies coefficientwise in $C^2\otimes\g$.  Applying Stokes
to the Lie-algebra-valued Bianchi identity yields
\begin{equation}
\label{eq:edge-boundary-bianchi}
 \epsilon c=\int_T P=\int_Y dP=-\int_Y[s,P].
\end{equation}
Thus $c$ is a linear expression in $P$ whose coefficients have positive
formal order.  Since
\[
 KP_0=-cK_{\mathrm{ext}}\rho,
 \qquad dP_0=-[s,P]-c\,d\rho,
\]
the contraction identity for $P_0$ gives the exact formula
\begin{equation}
\label{eq:edge-affine-curvature}
 P=i(r)-K[s,P]
       +c\bigl(\rho-dK_{\mathrm{ext}}\rho-Kd\rho\bigr).
\end{equation}
The original degree-three part of $K$ is defined on both $[s,P]$ and
$d\rho$, since all three-forms satisfy the boundary condition.
Substituting \eqref{eq:edge-boundary-bianchi} into
\eqref{eq:edge-affine-curvature} writes this as
\[
 P=i(r)+L_s(P),
\]
where $L_s$ raises $(h,\tau)$-adic order.  Hence
\[
 P=(1-L_s)^{-1}i(r).
\]
The residual components generate every curvature coefficient.  In the
other direction, $c=\epsilon^{-1}\int_TP$ and
$r=\operatorname{pr}(P-c\rho)$ are linear expressions in curvature
coefficients.  The residual and curvature ideals are therefore equal,
also over the nonreduced ring $S_\tau$.  This argument is valid for every
allowed extension $K_{\mathrm{ext}}$.

For a variation with $u$ fixed, the boundary term in
\eqref{eq:edge-variation} is
$-\epsilon\tr(u\dot v)/2$.  The added term in
\eqref{eq:edge-adjusted-potential} cancels it, so
\begin{equation}
\label{eq:edge-fixed-u-gradient}
 \frac{\partial f_\tau}{\partial h_j}
    =\left\langle\frac{\partial s(h,\tau)}{\partial h_j},P\right\rangle.
\end{equation}
As in the central case, this is a square matrix times $r$, with the
same perfect constant pairing.  The matrix is invertible over
$S_\tau$.  It follows that the internal Jacobian, residual, and full
curvature ideals coincide.  In particular the boundary commuting
equation $[u,v]=0$ is included in that ideal; it was not discarded.
\end{proof}

\subsection{Based gauge and marked holonomies}

The curvature and Jacobian ideals have been identified for both slices.
We now identify their quotient schemes with based representations and
track the two torus holonomies. Working on Artinian coefficient algebras
retains the nilpotent structure used by the ideal calculations.

We use the connection $d+a$ and parallel transport solving
$\dot U=-a(\dot\gamma)U$.  With the chosen boundary paths this gives
\begin{equation}
\label{eq:edge-holonomy-sign}
 G_\alpha=\exp(-u),\qquad G_\beta=\exp(-v)
\end{equation}
for a flat boundary connection $u\zeta+v\xi$.  Reversing the transport
convention changes the corresponding signs throughout.

Let $A$ be a local Artinian complex algebra with residue field
$\mathbb C$ and maximal ideal $\m_A$.  All connections and gauge
transformations in this paragraph reduce to the trivial ones modulo
$\m_A$.  The based gauge group consists of transformations equal to
the identity at the fixed point of $T$.  On a connection already in
the strip boundary form, we subsequently use gauge transformations
equal to the identity on all of $T$.

\begin{lemma}
\label{lem:edge-based-slice}
The central slice identifies its curvature-zero scheme with the formal
scheme of based representations of $\pi_1Y$ satisfying $G_\alpha=1$.
Over the dual numbers, the affine slice identifies its curvature-zero
scheme with the based representations satisfying
$G_\alpha=\exp(-\tau T_0)$.  These statements hold on all local
Artinian coefficient algebras, not only on reduced points.
\end{lemma}

\begin{proof}
Flat small connections modulo based gauge are equivalent to based
small representations by parallel transport.  This equivalence can
be constructed coefficientwise.  A representation defines the flat
bundle with its given framing; its small transition functions can be
trivialized smoothly by induction over the nilpotent filtration.
At each central square-zero stage this is the trivialization of an
additive cocycle of smooth functions, which follows from a partition
of unity.  Conversely, parallel transport is a finite iterated-integral
expression at each nilpotent order.  Flatness makes it invariant under
based homotopies.  Two flat connections with the same based monodromy
are related by the unique based gauge transformation obtained by
comparing their parallel transports.

Suppose first that a based representation has
$G_\alpha=\exp(-u)$, where $u=0$ or $u=\tau T_0$.  Put
$v=-\log G_\beta$.  The two boundary monodromies commute.  In the
affine case this implies $[u,v]=0$: conjugating the square-zero
exponential gives $(\exp(\ad_v)-1)u=0$, and
\[
 \exp(\ad_v)-1
    =\ad_v\left(1+\frac{\ad_v}{2!}
                   +\frac{(\ad_v)^2}{3!}+\cdots\right),
\]
whose factor in parentheses is invertible for nilpotent coefficients.
The assertion is automatic when $u=0$.  Hence
$u\zeta+v\xi$ is a flat torus connection with exactly the given based
monodromy.  Based parallel transport gives a gauge transformation
normalizing the original boundary connection to this strip form.
Its logarithm extends smoothly from the boundary, so the normalization
extends over $Y$.

Now use gauge transformations equal to the identity on $T$ to impose
\[
 K(a-u\widehat\zeta)=0.
\]
Their infinitesimal action is $a\mapsto a+d\phi$, with
$\phi\in C^0\otimes\g\otimes\m_A$, and $Kd\phi=\phi$.
Thus the exact slice coordinate can be killed uniquely, successively
in each nilpotent order.  For a flat connection in this gauge, setting
$h=\operatorname{pr}(a-u\widehat\zeta)$ and applying $K$ or
$K_{\mathrm{ext}}$ to its zero curvature gives exactly the corresponding
recursion.  Therefore the normalized connection is $s(h)$ or
$s(h,\tau)$.

There is no additional quotient by boundary conjugation.  If two
normalized boundary connections have the same based monodromy, their
coefficients $u,v$ agree by the formal logarithm.  A based gauge
transformation between those identical boundary connections is the
identity everywhere on $T$, again by parallel transport.  The subsequent
slice gauge is unique by $Kd=\id$ on $C^0$.  These constructions are
natural under maps of Artinian algebras, proving the assertions about
formal schemes, including their nilpotent structure.
\end{proof}

\Needspace{14\baselineskip}
We collect the conclusions, including the first-order coefficient
needed later.

\begin{proposition}[The marked edge potential]
\label{prop:edge-potential}
With the boundary polarization fixed above, there is a formal potential
$f\in\mathbb C[[h_1,\ldots,h_N]]$, where
$N=3(b_1(Y)-1)$, and a slice $s(h)$ such that:
\begin{enumerate}[label=\textup{(\roman*)}]
 \item $f=\mathcal C(s)$ has order at least three.  Its full Jacobian
 ideal equals the curvature ideal of $s$.  The formal critical scheme
 $\Crit(f)=\Spf(S/\Jac(f))$ is the based flat-representation scheme
 with $G_\alpha=1$, and carries its marked based matrix $G_\beta$.
 \item For each fixed $T_0\in\g$ there is an affine slice with boundary
 $\tau T_0\zeta+v(h,\tau)\xi$ and adjusted potential
 \[
  f_\tau=f+\tau b
    =\mathcal C(s(h,\tau))
             +\tfrac{\epsilon}{2}\tr(\tau T_0v(h,\tau)),
  \qquad \tau^2=0.
 \]
 Its internal Jacobian ideal equals its full curvature ideal.
 Accordingly $\Crit_h(f_\tau)$ is the based flat-representation scheme
 with $G_\alpha=\exp(-\tau T_0)$.  No derivative with respect to
 $\tau$ is imposed in this relative critical scheme.
 \item On the entire, possibly nonreduced, central critical scheme,
 \begin{equation}
 \label{eq:edge-log-coefficient}
  b=\epsilon\tr(T_0v(h,0))
      =-\epsilon\tr\bigl(T_0\log G_\beta\bigr)
       \quad\bmod\Jac(f).
 \end{equation}
 The function $b$ has order at least two.  Every flat connection with
 the indicated normalized boundary value is carried uniquely to the
 corresponding slice by gauge transformations equal to the identity
 on $T$.
\end{enumerate}
\end{proposition}

\begin{proof}
The preceding lemmas prove all assertions except the identification
of $b$ and its order.  For a general first variation of the two
boundary coefficients, the boundary part of
\eqref{eq:edge-variation} is
\[
 -\tfrac{\epsilon}{2}\tr(u\dot v-v\dot u).
\]
After adding the variation of $\epsilon\tr(uv)/2$, this becomes
$\epsilon\tr(v\dot u)$.  Thus the adjusted variation is
\begin{equation}
\label{eq:edge-adjusted-variation}
 \dot f_\tau
    =\int_Y\tr(\dot s\wedge P(s))
                      +\epsilon\tr(v\dot u).
\end{equation}
Take the coefficient derivative in $\tau$ at $\tau=0$, so that
$\dot u=T_0$.  On $S/\Jac(f)$ the central curvature vanishes as an
actual coefficientwise formal form, by
Lemma~\ref{lem:edge-central-ideal}.  The integral term therefore
vanishes in that quotient.  Equation~\eqref{eq:edge-holonomy-sign}
then proves \eqref{eq:edge-log-coefficient}.  This uses equality of
ideals and does not pass to the reduced critical locus.

The linear part of the affine slice is
$i(h)+\tau T_0\widehat\zeta$, which is closed and has no $\xi$
boundary coefficient.  Hence its Chern--Simons expression and its
boundary adjustment have joint $(h,\tau)$-order at least three.
The coefficient of $\tau$ consequently has $h$-order at least two.
All formal exponentials and logarithms used here are defined by their
nilpotent truncations on Artinian algebras and their formal limits.
\end{proof}
